% =========================================================
% Proof: Constant Multiple Rule
% Source: volume-iii/analysis/differentiation/notes/notes-algebra.tex
% Mode: standard
% =========================================================

\subsection*{Constant Multiple Rule}
\label{prf:constant-multiple-rule}

\begin{remark*}[Return]
\hyperref[thm:constant-multiple-rule]{$\leftarrow$ Back to Theorem (Constant Multiple Rule) in Notes}
\end{remark*}

\begin{proof}
Let $I \subseteq \mathbb{R}$ be an open interval, let $c \in I$, let
$f : I \to \mathbb{R}$ be differentiable at $c$, and let $\alpha \in \mathbb{R}$.
Set $h := \alpha f$. The claim is that $h$ is differentiable at $c$ and
$h'(c) = \alpha f'(c)$.

\ProofStep{1}{Form the difference quotient of $h$.}
For $x \in I$ with $x \neq c$, substituting $h(x) = \alpha f(x)$ gives
\[
  \frac{h(x) - h(c)}{x - c}
  = \frac{\alpha f(x) - \alpha f(c)}{x - c}
  = \alpha \cdot \frac{f(x) - f(c)}{x - c}.
\]

\ProofStep{2}{Pass to the limit.}
By \ByThm{Constant Multiple Rule for Limits},
\[
  \lim_{x \to c} \frac{h(x) - h(c)}{x - c}
  = \alpha \cdot \lim_{x \to c} \frac{f(x) - f(c)}{x - c}.
\]
Since $f$ is differentiable at $c$, the limit on the right exists and equals
$f'(c)$ by \ByDef{Differentiability at a Point}. Therefore $h'(c) = \alpha f'(c)$.
\end{proof}

\begin{remark*}[Proof shape]
The scalar $\alpha$ factors out of the difference quotient by arithmetic,
reducing the problem immediately to the known limit $f'(c)$. The Constant
Multiple Rule for Limits handles the passage through the scalar.

\FlashProof{1. Factor $\alpha$ out of the difference quotient of $\alpha f$.
  2. Apply the Constant Multiple Rule for Limits; differentiability of $f$
  supplies $f'(c)$.}
\end{remark*}

\begin{remark*}[Dependencies]
\begin{itemize}
  \item \hyperref[def:derivative-at-point]{Definition of Differentiability at a Point}:
    Identifies the limit of the difference quotient of $f$ as $f'(c)$.
  \item \hyperref[thm:limit-constant-multiple]{Constant Multiple Rule for Limits}:
    Licenses pulling $\alpha$ outside the limit.
\end{itemize}

\FlashUses{def:derivative-at-point, thm:limit-constant-multiple}
\end{remark*}